\documentclass[a4paper]{article}

\usepackage[a4paper, top=8mm, bottom=8mm, left=8mm, right=8mm]{geometry}

\usepackage{polyglossia}
\setdefaultlanguage[babelshorthands=true]{russian}
\setotherlanguage{english}

\usepackage{fontspec}
\setmainfont{FreeSerif}
\newfontfamily{\russianfonttt}[Scale=0.7]{DejaVuSansMono}

\usepackage[tiny, compact]{titlesec}

\usepackage{titling}
\setlength{\droptitle}{-1cm}
\pretitle{\begin{center}\begin{bfseries}\Large}
\posttitle{\par\end{bfseries}\end{center}}
\preauthor{\begin{center}\normalsize}
\postauthor{\par\end{center}\vspace{-1.8cm}}

\usepackage{hyperref}
\usepackage{bookmark}
\usepackage{csquotes}
\usepackage{amsmath}
\usepackage{booktabs}
\usepackage{array}

\setlength{\parindent}{1.25em}
\setlength{\parskip}{0pt}

\titleformat{\section}
{\normalfont\bfseries\large}
{\thesection.}
{0.5em}
{}

\titleformat{\subsection}
{\normalfont\bfseries\normalsize}
{\thesubsection.}
{0.5em}
{}

\title{Исследование способов оптимизации операции перемножения плотных матриц (SGEMM) на архитектуре AMD RDNA3 с использованием OpenCL}

\author{Гарин Александр Евгеньевич}

\date{}

\hypersetup{
    pdftitle={Исследование способов оптимизации операции перемножения плотных матриц (SGEMM) на архитектуре AMD RDNA3 с использованием OpenCL},
    pdfauthor={Гарин Александр Евгеньевич},
    hidelinks
}

\begin{document}

\maketitle

\begin{flushright}
    Группа: \emph{25.Б71-мм}

    Кафедра: \emph{кафедра системного программирования}

    Научный руководитель: \emph{Григорьев Семен Вячеславович}

    Номер семестра практики: \emph{2}
\end{flushright}

\section{Введение}

SGEMM (Single-precision General Matrix Multiplication)~--- операция
перемножения плотных матриц, элементы которых являются числами с
плавающей точкой одинарной точности. Эта операция является одной из
базовых в вычислительной линейной алгебре, и от скорости её выполнения
зависит скорость многих прикладных вычислений. Для ускорения SGEMM
часто используют графические процессоры (GPU, Graphics Processing
Unit), которые выполняют одновременно большое число однотипных операций.

Работа выполнялась на основе учебного материала Cedric
Nugteren\footnote{URL: \url{https://cnugteren.github.io/tutorial/pages/page1.html}
(дата обращения: 05.09.2026).}, в котором реализация SGEMM
последовательно улучшается за десять шагов: каждый шаг вводит новый
способ организации вычислений или доступа к данным. Программы в этом
материале написаны с использованием OpenCL (Open Computing
Language)\footnote{URL: \url{https://www.khronos.org/opencl/}
(дата обращения: 19.09.2026).}~--- открытого стандарта для написания
программ, выполняемых на графических процессорах и других
вычислительных устройствах разных производителей.

Результаты в учебном материале получены на дискретном графическом
процессоре NVIDIA Tesla K40m. Эффект от одних и тех же приёмов
оптимизации и наилучшие значения их параметров зависят от устройства,
на котором выполняются вычисления, поэтому результаты, полученные на
одном графическом процессоре, нельзя без проверки переносить на другой.
В данной работе приёмы из учебного материала исследуются на
интегрированном графическом процессоре AMD Radeon 760M, то есть
графическом процессоре, расположенном на одном кристалле с центральным
процессором. Radeon 760M построен на
архитектуре RDNA3\footnote{URL: \url{https://www.amd.com/en/technologies/rdna.html}
(дата обращения: 13.09.2026).}~--- семействе архитектур графических
процессоров AMD. Интегрированные графические процессоры AMD
устанавливаются во многие ноутбуки, поэтому полученные данные могут
быть полезны при выборе варианта реализации и параметров SGEMM на
OpenCL для подобных устройств.

\section{Постановка задачи}

Целью работы является сравнение производительности вариантов
реализации SGEMM, построенных на основе шагов учебного материала, на
графическом процессоре AMD Radeon 760M при параметрах, заданных по
умолчанию, и при параметрах, подобранных перебором всех допустимых
значений.

Для достижения цели были поставлены следующие задачи:
\begin{itemize}
    \item реализовать на основе шагов учебного материала варианты SGEMM
    и запустить их на используемой программно-аппаратной платформе;
    \item измерить время выполнения и производительность вариантов при
    параметрах, заданных по умолчанию, и сравнить их между собой;
    \item определить допустимые значения параметров исходя из
    ограничений устройства и разработать инструмент, перебирающий все
    такие значения;
    \item сравнить варианты с подобранными параметрами;
    \item оценить снижение производительности, вызванное поддержкой
    матриц произвольного размера;
    \item настроить автоматическую проверку корректности всех вариантов
    при изменении кода.
\end{itemize}

\section{Описание решения}

\subsection{Основные понятия OpenCL}

Программа на OpenCL состоит из двух частей. Основная программа
выполняется на центральном процессоре, подготавливает данные и
запускает вычисления. Ядро (kernel)~--- функция, которая выполняется
на графическом процессоре одновременно большим числом потоков; каждый
поток обрабатывает свою часть данных. Потоки объединяются в рабочие
группы~--- группы потоков, которые выполняются совместно и могут
обмениваться данными через общую память. Запуск ядра выполняется через
очередь команд: основная программа ставит в неё команду запуска и может
ожидать её завершения.

В работе используются три вида памяти графического процессора:
\begin{itemize}
    \item глобальная память~--- память, доступная всем потокам, в
    которой хранятся исходные матрицы и матрица результата;
    у интегрированного графического процессора она располагается в
    оперативной памяти компьютера;
    \item локальная память~--- небольшая быстрая память, общая для
    потоков одной рабочей группы; в неё загружают данные, к которым
    потоки группы обращаются многократно;
    \item регистры~--- самая быстрая память, доступная только одному
    потоку.
\end{itemize}

Все матрицы хранятся в памяти по столбцам (column-major), как и в
учебном материале: элемент, расположенный в строке $r$ и столбце $c$
матрицы из $R$ строк, имеет в одномерном массиве индекс $i = r + cR$.
Вычисляется произведение $C = AB$, где $A$ имеет размер $M \times K$,
$B$~--- $K \times N$, а $C$~--- $M \times N$. Далее размеры задачи
записываются в виде $M \times N \times K$.

\subsection{Реализованные варианты}

Номера вариантов совпадают с номерами шагов учебного материала,
на основе которых они написаны. Варианты являются упрощёнными
реализациями идей соответствующих шагов и в ряде случаев отличаются
от кода учебного материала; отличия указаны ниже. Реализованы следующие
варианты:
\begin{itemize}
    \item вариант~1 (базовый)~--- каждый поток вычисляет один элемент
    матрицы $C$, считывая строку матрицы $A$ и столбец матрицы $B$
    непосредственно из глобальной памяти;
    \item вариант~2~--- добавлено блокирование: матрицы разбиваются на
    квадратные блоки размера $TS \times TS$, рабочая группа загружает
    очередные блоки матриц $A$ и $B$ в локальную память, после чего
    потоки выполняют вычисления, обращаясь только к локальной памяти;
    \item вариант~3~--- на основе варианта~2 каждый поток вычисляет
    не один, а $WPT$ элементов одной строки матрицы $C$, храня
    промежуточные суммы в регистрах; одно считанное из локальной памяти
    значение блока матрицы $A$ используется для вычисления $WPT$
    элементов результата;
    \item вариант~4~--- на основе варианта~2 данные загружаются,
    обрабатываются и сохраняются с помощью векторного типа OpenCL
    \texttt{float4}, содержащего четыре числа, поэтому каждый поток
    вычисляет четыре соседних элемента столбца матрицы $C$;
    \item вариант~5~--- на основе варианта~3 блок матрицы $B$ хранится
    в локальной памяти в транспонированном виде, а строка этого блока
    удлинена на один элемент; такой сдвиг применяется для уменьшения
    конфликтов банков локальной памяти, то есть ситуаций, когда
    несколько потоков одновременно обращаются к одному и тому же
    участку (банку) локальной памяти и их обращения выполняются по
    очереди; в учебном материале на этом шаге вместо этого
    транспонируется вся матрица в глобальной памяти и используются
    прямоугольные блоки;
    \item вариант~6~--- двумерное блокирование в регистрах: каждый поток
    вычисляет фрагмент матрицы $C$ из $WPT \times WPT$ элементов,
    расположенных с шагом $TS/WPT$ по строкам и столбцам;
    \item вариант~7~--- на основе варианта~6 на каждом шаге вычислений
    поток один раз считывает нужные значения из локальной памяти в
    отдельные переменные в регистрах, после чего выполняет с ними все
    умножения и сложения; в учебном материале на этом шаге вместо этого
    используется загрузка данных векторными типами;
    \item вариант~10~--- на основе варианта~3 добавлена обработка
    матриц, размеры которых не кратны размеру блока: число блоков
    округляется вверх, элементы за границами матриц при загрузке
    заменяются нулями, а запись в матрицу $C$ выполняется только в
    пределах её размеров; в учебном материале этот шаг построен на
    основе более сложного варианта с двумерным блокированием.
\end{itemize}

Шаги~8 и~9 учебного материала в работу не вошли. Шаг~8 посвящён
оптимизациям, специфичным для CUDA\footnote{URL: \url{https://developer.nvidia.com/cuda-toolkit}
(дата обращения: 19.09.2026).}~--- платформы параллельных вычислений
компании NVIDIA,~--- и архитектуры графических процессоров NVIDIA
Kepler, поэтому он неприменим к используемому устройству. При попытке
запуска варианта, соответствующего шагу~9 (предварительная загрузка
следующих блоков данных во время вычислений), система переставала
отвечать: изображение на экране замирало и не реагировало на действия
пользователя. Причина такого поведения в работе не исследовалась,
поэтому измерения для шага~9 не проводились.

\subsection{Параметры вариантов и их допустимые значения}
\label{sec:params}

Варианты используют следующие параметры:
\begin{itemize}
    \item \texttt{TS}~--- размер квадратного блока;
    \item \texttt{WPT} (Work Per Thread)~--- количество элементов
    результата, вычисляемых одним потоком (в вариантах~6 и~7~---
    по каждому из двух измерений);
    \item \texttt{WIDTH}~--- число элементов в векторном типе
    (только в варианте~4, где оно равно~4 и не изменялось, поскольку
    код ядра написан для типа \texttt{float4}).
\end{itemize}

Размер рабочей группы определяется параметрами: в вариантах~1 и~2 он
равен $TS^2$ потоков, в вариантах~3, 5 и~10~--- $TS^2/WPT$,
в варианте~4~--- $TS^2/4$, в вариантах~6 и~7~--- $(TS/WPT)^2$.
Варианты~2--7 и~10 хранят в локальной памяти два блока из $TS^2$ чисел,
то есть используют около $8 TS^2$ байт локальной памяти; вариант~1
локальную память не использует.

Допустимые значения параметров определяются следующими условиями:
\begin{itemize}
    \item значение $TS$ должно делить размер задачи $4096 = 2^{12}$
    нацело, поскольку варианты~1--7 не обрабатывают неполные блоки,
    поэтому $TS$ может быть только степенью двойки;
    \item значение $WPT$ должно делить $TS$ нацело, поэтому $WPT$ также
    является степенью двойки;
    \item размер рабочей группы должен быть не больше 256 потоков~---
    максимального размера рабочей группы, который сообщает устройство;
    \item размер рабочей группы должен быть не меньше 32 потоков:
    графический процессор выполняет потоки пачками (волнами, wavefront)
    по 32 потока, и группа меньшего размера занимает волну не
    полностью;
    \item объём используемой локальной памяти не должен превышать
    64~КБ~--- объёма локальной памяти, доступной рабочей группе, поэтому
    $TS \le 64$.
\end{itemize}
Значения ограничений устройства получены с помощью утилиты
\texttt{clinfo}\footnote{URL: \url{https://github.com/Oblomov/clinfo}
(дата обращения: 19.09.2026).}, выводящей характеристики устройств
OpenCL. При переборе проверялись все комбинации параметров,
удовлетворяющие этим условиям (табл.~\ref{tab:params}). Например,
для вариантов~1 и~2 допустимы только $TS=8$ и $TS=16$: при $TS=4$
рабочая группа содержит 16 потоков, а при $TS=32$~--- 1024 потока.

Параметрами по умолчанию в работе называются значения, с которыми
варианты были реализованы до проведения перебора. В учебном материале
для вариантов~1--5 используется \texttt{TS=32}, однако при этом
значении вариантам~1 и~2 потребовалась бы рабочая группа из 1024
потоков, поэтому для вариантов~1--5 и~10 по умолчанию выбрано значение
\texttt{TS=16}. Для вариантов~6 и~7 значения \texttt{TS=32},
\texttt{WPT=4} выбраны автором при реализации; в учебном материале
для этих шагов используются блоки размера $128 \times 128$ и фрагменты
из $8 \times 8$ элементов на поток.

\begin{table}[h!]
\centering
\caption{Параметры по умолчанию и допустимые значения, проверенные при переборе}
\label{tab:params}
\begin{tabular}{c l l c}
\toprule
Вариант & По умолчанию & Допустимые значения & Число комбинаций \\
\midrule
1, 2      & \texttt{TS=16}                 & $TS \in \{8, 16\}$                                 & 2 \\
3, 5, 10  & \texttt{TS=16}, \texttt{WPT=8} & $TS \in \{8, 16, 32, 64\}$, $32 \le TS^2/WPT \le 256$ & 13 \\
4         & \texttt{TS=16}                 & $TS \in \{16, 32\}$                                & 2 \\
6         & \texttt{TS=32}, \texttt{WPT=4} & $TS \in \{8, 16, 32, 64\}$, $TS/WPT \in \{8, 16\}$  & 7 \\
7         & \texttt{TS=32}, \texttt{WPT=4} & $TS \in \{32, 64\}$, $WPT = 4$                     & 2 \\
\bottomrule
\end{tabular}
\end{table}

\subsection{Программа измерений и инструмент перебора}

Каждый вариант сначала был реализован автором в виде самостоятельной
программы на C++, построенной на основе минимального примера из
репозитория учебного материала\footnote{URL: \url{https://github.com/CNugteren/myGEMM}
(дата обращения: 19.09.2026).}: варианты~1--3, 5--7 и~10~--- в
последовательных версиях файла \texttt{extra/minimal.cpp}, вариант~4~---
в файле \texttt{extra/gemm4\_vector.cpp}. Для перебора параметров тексты
ядер всех вариантов без изменений собраны автором в одной программе
\texttt{extra/tune/tune\_gemm.cpp}. Номер варианта, значения
параметров и размер задачи задаются при компиляции программы;
значения параметров передаются в ядро при его компиляции. Программа
инициализирует матрицы, запускает ядро, измеряет время выполнения,
вычисляет производительность и разброс, проверяет корректность
результата и выводит все значения одной строкой.

Скрипт \texttt{extra/tune/tune\_all.sh}, также разработанный автором,
сам определяет все комбинации параметров, удовлетворяющие условиям из
подраздела~\ref{sec:params}, и для каждой из них компилирует программу,
запускает её и дописывает строку результата в файл
\texttt{tune\_results.txt} сразу после её получения, поэтому при
аварийном завершении уже полученные результаты сохраняются. Всего
проверяется 67 комбинаций: 41 для вариантов~1--7 и по 13 для
варианта~10 на каждом из двух размеров задачи. Исходный код и файл
результатов размещены в репозитории
проекта\footnote{URL: \url{https://github.com/aleksandrgarin/myGEMM/tree/practice-optimization}
(дата обращения: 22.09.2026).}.

\subsection{Автоматическая проверка корректности}

В репозитории настроена непрерывная интеграция (CI, Continuous
Integration)~--- автоматическая сборка и проверка программы при каждом
изменении кода~--- с помощью GitHub
Actions\footnote{URL: \url{https://docs.github.com/en/actions}
(дата обращения: 22.09.2026).}. При каждой отправке изменений в
репозиторий на сервере GitHub устанавливается PoCL (Portable Computing
Language)\footnote{URL: \url{https://portablecl.org/}
(дата обращения: 22.09.2026).}~--- реализация OpenCL, выполняющая ядра
на центральном процессоре,~--- после чего скрипт перебора собирает и
запускает все 67 комбинаций параметров на матрицах размера
$256 \times 256$ (для варианта~10 с неполными блоками~--- размер задачи
$257 \times 160 \times 256$). Проверка завершается с ошибкой, если хотя
бы одна комбинация не собралась, не выполнилась или дала неверный
результат. Производительность в CI не измеряется, поскольку серверы
GitHub Actions не оснащены графическим процессором.

\section{Эксперименты}

\subsection{Окружение}

Эксперименты проводились на ноутбуке HONOR MagicBook X 16 (модель
BORN-H5651) с процессором AMD Ryzen 5 7640HS и интегрированным
графическим процессором AMD Radeon 760M. Программно-аппаратная
конфигурация:
\begin{itemize}
    \item архитектура графического процессора~--- RDNA3, идентификатор
    устройства~--- \texttt{gfx1103};
    \item операционная система~--- Ubuntu 24.04.4 LTS, версия ядра
    операционной системы Linux~--- 6.17.0-1020-oem;
    \item компилятор GCC/G++\footnote{URL: \url{https://gcc.gnu.org/onlinedocs/13.3.0/}
    (дата обращения: 13.09.2026).}~--- версия 13.3.0;
    \item ROCm (Radeon Open Compute)\footnote{URL: \url{https://rocm.docs.amd.com/}
    (дата обращения: 13.09.2026).}~--- программная платформа AMD для
    вычислений на графических процессорах, версия 6.4.4;
    \item среда выполнения OpenCL~--- пакет
    \texttt{rocm-opencl-runtime}\footnote{URL: \url{https://rocm.docs.amd.com/projects/radeon-ryzen/en/latest/docs/install/installryz/native_linux/install-package-manager.html}
    (дата обращения: 13.09.2026).} версии 6.4.4;
    \item платформа OpenCL~--- AMD Accelerated Parallel Processing
    (AMD-APP), версия 3649.0, поддерживает спецификацию OpenCL~2.1;
    \item версия OpenCL, поддерживаемая устройством,~--- 2.0;
    \item версия драйвера~--- 3649.0 (HSA1.1, LC), где HSA
    (Heterogeneous System Architecture)~--- архитектура взаимодействия
    центрального и графического процессоров, а LC~--- режим компиляции
    ядер;
    \item загрузчик OpenCL ICD (Installable Client
    Driver)\footnote{URL: \url{https://github.com/KhronosGroup/OpenCL-ICD-Loader}
    (дата обращения: 13.09.2026).}~--- компонент, обеспечивающий
    обнаружение и выбор реализаций OpenCL, версия 2.3.2;
    \item максимальный размер рабочей группы~--- 256 потоков, ширина
    волны~--- 32 потока, объём локальной памяти~--- 64~КБ (по данным
    \texttt{clinfo}).
\end{itemize}

Во время измерений ноутбук работал от сети, в Ubuntu был выбран профиль
энергопотребления <<Производительность>>. Программы компилировались
по стандарту C++11 с параметрами \texttt{-O3} (высокий уровень
оптимизации) и \texttt{-Wall} (вывод диагностических предупреждений).

\subsection{Входные данные}

Для вариантов~1--7 использовались квадратные матрицы размера
$4096 \times 4096$ (размер задачи $4096 \times 4096 \times 4096$), как
в учебном материале. Этот размер кратен всем допустимым размерам
блока (8, 16, 32 и 64), поэтому неполные блоки не возникают.

Вариант~10 запускался с двумя размерами задачи. Размер
$4097 \times 4000 \times 4096$ приводит к неполным блокам: значение
4097 не делится ни на один из допустимых размеров блока, поэтому
неполные блоки по измерению $M$ возникают при всех значениях $TS$;
значение 4000 делится на 8, 16 и~32, но не на~64, поэтому при $TS=64$
неполные блоки возникают также по измерению $N$. Размер
$4096 \times 4096 \times 4096$ совпадает с размером задачи для
остальных вариантов и нужен, чтобы отделить влияние проверок границ от
влияния размеров матриц.

Элементы матриц инициализировались по формулам из учебного материала:
$$
A_i = 3.6i + i^2 + 3.1,
\qquad
B_i = 1.2i + 0.01i^2 + 13.9,
$$
где $i$~--- индекс элемента в одномерном массиве. В коде, как и в
учебном материале, величина $i^2$ для матрицы $A$ вычисляется в
32-битном целочисленном типе и при $i > 46340$ переполняется, поэтому
фактические значения элементов $A$ отличаются от приведённой формулы.
Объём вычислений от значений элементов не зависит. Матрица $C$ перед
вычислениями заполнялась нулями.

\subsection{Методика измерений}

Для каждой комбинации параметров ядро запускалось 10 раз подряд, как
в минимальном примере учебного материала, без предварительного
(прогревочного) запуска. Время каждого запуска измерялось на
центральном процессоре таймером \texttt{std::chrono} от постановки
команды запуска ядра в очередь до завершения её выполнения.
Инициализация матриц, компиляция ядра, выделение памяти и передача
данных между оперативной памятью и графическим процессором в измеряемый
интервал не входили.

По результатам запусков вычислялись среднее время
$$
t_{\mathrm{avg}} = \frac{1}{10}\sum_{i=1}^{10} t_i,
$$
где $t_i$~--- время $i$-го запуска в секундах, и производительность
$$
G = \frac{2MNK}{t_{\mathrm{avg}}\cdot 10^9},
$$
где $M$, $N$ и $K$~--- фактические размеры задачи. Множитель~2
учитывает одно умножение и одно сложение на каждую из $MNK$
элементарных операций. Производительность выражается в GFLOPS
(Giga Floating Point Operations Per Second)~--- миллиардах операций с
плавающей точкой в секунду.

Разброс характеризуется половиной размаха производительности:
$$
\Delta G = \frac{G_{\max}-G_{\min}}{2},
$$
где $G_{\max}$ и $G_{\min}$~--- производительность, вычисленная по
времени самого быстрого и самого медленного из десяти запусков.
Результаты записываются в виде $G \pm \Delta G$. Два результата
считаются различимыми, если интервалы $[G - \Delta G,\ G + \Delta G]$
не пересекаются. Лучшей для варианта называется комбинация параметров
с наибольшим значением $G$ среди проверенных.

Корректность вычислений проверялась для каждой комбинации параметров:
для 100 случайно выбранных элементов матрицы $C$ результат сравнивался
с вычислением на центральном процессоре с двойной точностью. Полная
проверка не проводилась, поскольку она потребовала бы около
$6.9 \cdot 10^{10}$ операций умножения и сложения на центральном
процессоре для каждой комбинации. Погрешность элемента вычислялась как
$\left|C_{rc} - C^{\mathrm{ref}}_{rc}\right| / \sum_k \left|A_{rk} B_{kc}\right|$.
Во всех 67 комбинациях ни один проверенный элемент не превысил
допустимую погрешность $10^{-3}$; максимальная погрешность составила
$8.3 \cdot 10^{-8}$ для размера задачи $4096 \times 4096 \times 4096$
и $1.1 \cdot 10^{-7}$ для размера $4097 \times 4000 \times 4096$, что
соответствует точности вычислений с числами одинарной точности.

Все приведённые результаты получены за один запуск скрипта перебора
22.09.2026 длительностью около 7~минут, то есть в одинаковых условиях.

\subsection{Сравнение вариантов при параметрах по умолчанию}
\label{sec:default}

Цель эксперимента~--- сравнить варианты~1--7 при параметрах по
умолчанию (табл.~\ref{tab:params}) на одинаковых матрицах
$4096 \times 4096$. Вариант~10 в это сравнение не включён, поскольку
его назначение и условия запуска отличаются; он рассматривается в
подразделе~\ref{sec:v10}. Результаты приведены в
табл.~\ref{tab:default}.

\begin{table}[h!]
\centering
\caption{Производительность вариантов SGEMM при параметрах по умолчанию}
\label{tab:default}
\begin{tabular}{c >{\raggedright\arraybackslash}p{7.2cm} l r r}
\toprule
Вариант & Описание & Параметры & Время, с & Производительность, GFLOPS \\
\midrule
1 & Базовый вариант                                                & \texttt{TS=16}                 & 2.626 & $52.3 \pm 0.5$ \\
2 & Блокирование в локальной памяти                                & \texttt{TS=16}                 & 0.476 & $288.5 \pm 16.7$ \\
3 & Несколько элементов результата на поток                        & \texttt{TS=16}, \texttt{WPT=8} & 0.288 & $476.5 \pm 34.2$ \\
4 & Векторный тип данных \texttt{float4}                           & \texttt{TS=16}                 & 0.277 & $496.5 \pm 68.2$ \\
5 & Транспонированное хранение блока в локальной памяти со сдвигом & \texttt{TS=16}, \texttt{WPT=8} & 0.440 & $312.4 \pm 17.1$ \\
6 & Двумерное блокирование в регистрах                             & \texttt{TS=32}, \texttt{WPT=4} & 0.318 & $432.7 \pm 38.9$ \\
7 & Двумерное блокирование с явной загрузкой в регистры            & \texttt{TS=32}, \texttt{WPT=4} & 0.305 & $450.9 \pm 41.7$ \\
\bottomrule
\end{tabular}
\end{table}

Блокирование в локальной памяти (вариант~2) повысило среднюю
производительность относительно базового варианта примерно в 5.5 раза,
а вычисление нескольких элементов на поток (вариант~3)~--- ещё
примерно в 1.7 раза. Наибольшую среднюю производительность при
параметрах по умолчанию показал вариант~4~--- 496.5~GFLOPS, что
примерно в 9.5 раза выше, чем у базового варианта. Однако результаты
вариантов~3, 4, 6 и~7 не различимы: их интервалы разброса пересекаются.
В частности, использование векторного типа \texttt{float4} (вариант~4)
не дало различимого прироста относительно варианта~3. Вариант~5
оказался медленнее варианта~3, на котором он основан, примерно на
34.4~\%, и это различие выходит за пределы разброса.

\subsection{Подбор параметров для всех вариантов}
\label{sec:tuning}

Результаты предыдущего эксперимента зависят от значений параметров по
умолчанию, которые для разных вариантов выбирались по-разному.
Цель данного эксперимента~--- проверить, сохранится ли соотношение
вариантов, если для каждого из них выбрать лучшую из всех допустимых
комбинаций параметров. Результаты перебора для вариантов~3 и~5,
у которых одинаковый набор допустимых комбинаций, приведены в
табл.~\ref{tab:tune35}, для остальных вариантов~--- в
табл.~\ref{tab:tuneother}.

\begin{table}[h!]
\centering
\caption{Результаты перебора параметров для вариантов 3 и 5}
\label{tab:tune35}
\begin{tabular}{c c c r r}
\toprule
 & & & \multicolumn{2}{c}{Производительность, GFLOPS} \\
\cmidrule(lr){4-5}
$TS$ & $WPT$ & Потоков в рабочей группе & Вариант~3 & Вариант~5 \\
\midrule
8  & 1  & 64  & $228.0 \pm 9.3$  & $226.0 \pm 9.5$  \\
8  & 2  & 32  & $280.5 \pm 13.6$ & $279.1 \pm 12.4$ \\
16 & 1  & 256 & $288.9 \pm 17.4$ & $287.5 \pm 17.5$ \\
16 & 2  & 128 & $401.0 \pm 33.8$ & $389.1 \pm 31.6$ \\
16 & 4  & 64  & $500.5 \pm 68.1$ & $508.0 \pm 60.5$ \\
16 & 8  & 32  & $476.5 \pm 34.2$ & $312.4 \pm 17.1$ \\
32 & 4  & 256 & $486.8 \pm 45.5$ & $523.2 \pm 55.3$ \\
32 & 8  & 128 & $455.1 \pm 38.9$ & $438.6 \pm 36.2$ \\
32 & 16 & 64  & $448.5 \pm 36.4$ & $334.0 \pm 18.7$ \\
32 & 32 & 32  & $190.1 \pm 7.8$  & $169.7 \pm 5.7$  \\
64 & 16 & 256 & $629.7 \pm 78.0$ & $602.3 \pm 69.5$ \\
64 & 32 & 128 & $485.0 \pm 47.2$ & $610.6 \pm 75.6$ \\
64 & 64 & 64  & $142.7 \pm 6.0$  & $443.7 \pm 37.2$ \\
\bottomrule
\end{tabular}
\end{table}

\begin{table}[h!]
\centering
\caption{Результаты перебора параметров для вариантов 1, 2, 4, 6 и 7}
\label{tab:tuneother}
\begin{tabular}{c c c c r}
\toprule
Вариант & $TS$ & $WPT$ & Потоков в рабочей группе & Производительность, GFLOPS \\
\midrule
1 & 8  & --- & 64  & $52.0 \pm 0.5$ \\
1 & 16 & --- & 256 & $52.3 \pm 0.5$ \\
\midrule
2 & 8  & --- & 64  & $226.8 \pm 9.6$ \\
2 & 16 & --- & 256 & $288.5 \pm 16.7$ \\
\midrule
4 & 16 & --- & 64  & $496.5 \pm 68.2$ \\
4 & 32 & --- & 256 & $419.3 \pm 36.4$ \\
\midrule
6 & 8  & 1 & 64  & $229.4 \pm 9.8$ \\
6 & 16 & 1 & 256 & $289.6 \pm 17.0$ \\
6 & 16 & 2 & 64  & $446.4 \pm 71.6$ \\
6 & 32 & 2 & 256 & $427.8 \pm 38.1$ \\
6 & 32 & 4 & 64  & $432.7 \pm 38.9$ \\
6 & 64 & 4 & 256 & $819.6 \pm 123.2$ \\
6 & 64 & 8 & 64  & $683.9 \pm 87.3$ \\
\midrule
7 & 32 & 4 & 64  & $450.9 \pm 41.7$ \\
7 & 64 & 4 & 256 & $820.8 \pm 123.3$ \\
\bottomrule
\end{tabular}
\end{table}

Варианты~3 и~6 при $WPT=1$ выполняют те же вычисления, что и
вариант~2, и их результаты совпадают с результатами варианта~2 при том
же $TS$ в пределах разброса (например, 228.0, 229.4 и 226.8~GFLOPS при
$TS=8$). Это подтверждает, что различия между вариантами вызваны
изменениями в ядрах, а не условиями измерения.

У варианта~3 производительность растёт при увеличении $WPT$ до
некоторого значения, а при $WPT = TS$ резко падает (190.1~GFLOPS при
$TS=32$ и 142.7~GFLOPS при $TS=64$). Лучшей оказалась комбинация
$TS=64$, $WPT=16$~--- 629.7~GFLOPS, что различимо выше, чем при
параметрах по умолчанию (476.5~GFLOPS). Влияние транспонированного
хранения со сдвигом (вариант~5) зависит от параметров: в 8 из 13
комбинаций результаты вариантов~3 и~5 не различимы, в трёх
комбинациях ($TS=16$, $WPT=8$; $TS=32$, $WPT=16$; $TS=32$, $WPT=32$)
вариант~5 различимо медленнее, а в двух ($TS=64$, $WPT=32$ и
$TS=64$, $WPT=64$)~--- различимо быстрее. Лучшие результаты
вариантов~3 и~5 не различимы. У варианта~1 значения $TS$ не дают
различимой разницы, у варианта~2 значение $TS=16$ различимо лучше
$TS=8$, у варианта~4 результаты двух допустимых значений $TS$ не
различимы. У вариантов~6 и~7 наибольшая производительность получена
при $TS=64$, $WPT=4$, а результаты вариантов~6 и~7 при одинаковых
параметрах совпадают в пределах разброса, то есть явная загрузка
значений в регистры (вариант~7) не дала различимого эффекта.

Для сравнения вариантов между собой в табл.~\ref{tab:best} для каждого
из них приведена лучшая комбинация из табл.~\ref{tab:tune35}
и~\ref{tab:tuneother}.

\begin{table}[h!]
\centering
\caption{Производительность вариантов SGEMM при лучших из допустимых параметров}
\label{tab:best}
\begin{tabular}{c l r r}
\toprule
Вариант & Лучшие параметры & Время, с & Производительность, GFLOPS \\
\midrule
1 & \texttt{TS=16}                  & 2.626 & $52.3 \pm 0.5$ \\
2 & \texttt{TS=16}                  & 0.476 & $288.5 \pm 16.7$ \\
3 & \texttt{TS=64}, \texttt{WPT=16} & 0.218 & $629.7 \pm 78.0$ \\
4 & \texttt{TS=16}                  & 0.277 & $496.5 \pm 68.2$ \\
5 & \texttt{TS=64}, \texttt{WPT=32} & 0.225 & $610.6 \pm 75.6$ \\
6 & \texttt{TS=64}, \texttt{WPT=4}  & 0.168 & $819.6 \pm 123.2$ \\
7 & \texttt{TS=64}, \texttt{WPT=4}  & 0.167 & $820.8 \pm 123.3$ \\
\bottomrule
\end{tabular}
\end{table}

После подбора параметров соотношение вариантов изменилось.
Наибольшую среднюю производительность показали варианты~6 и~7 с
параметрами $TS=64$, $WPT=4$~--- 819.6 и 820.8~GFLOPS соответственно;
между собой они не различимы. Нижние границы их интервалов разброса
(696.4 и 697.5~GFLOPS) выше верхних границ интервалов вариантов~1, 2,
4 и~5 (наибольшая из них~--- 686.2~GFLOPS у варианта~5), поэтому их
превосходство над этими вариантами выходит за пределы разброса.
С вариантом~3 при $TS=64$, $WPT=16$ (629.7~GFLOPS, верхняя граница
интервала 707.7~GFLOPS) интервалы пересекаются, поэтому превосходство
вариантов~6 и~7 над ним по полученным данным не подтверждается, хотя
их средняя производительность примерно в 1.3 раза выше. Средняя
производительность вариантов~6 и~7 примерно в 15.7 раза выше, чем у
базового варианта. Вариант с наибольшей средней производительностью
при параметрах по умолчанию (вариант~4) после подбора параметров
различимо уступает вариантам~6 и~7: для него допустимы только
$TS \le 32$, а у вариантов~3, 5, 6 и~7 лучшие результаты получены
при $TS=64$.

\subsection{Затраты на поддержку матриц произвольного размера}
\label{sec:v10}

Вариант~10 отличается от варианта~3 только проверками границ, которые
позволяют обрабатывать матрицы произвольного размера. Цель
эксперимента~--- определить, насколько эти проверки снижают
производительность и влияет ли на неё наличие неполных блоков. Для
этого вариант~10 запускался со всеми допустимыми комбинациями
параметров для обоих размеров задачи, а его результаты сравнивались с
результатами варианта~3 при тех же параметрах. Результаты приведены в
табл.~\ref{tab:v10}.

\begin{table}[h!]
\centering
\caption{Производительность вариантов 3 и 10 при одинаковых параметрах}
\label{tab:v10}
\begin{tabular}{c c r r r}
\toprule
 & & \multicolumn{3}{c}{Производительность, GFLOPS} \\
\cmidrule(lr){3-5}
$TS$ & $WPT$ & Вариант~3, & Вариант~10, & Вариант~10, \\
 & & $4096 \times 4096 \times 4096$ & $4096 \times 4096 \times 4096$ & $4097 \times 4000 \times 4096$ \\
\midrule
8  & 1  & $228.0 \pm 9.3$  & $151.7 \pm 4.4$  & $159.8 \pm 4.4$  \\
8  & 2  & $280.5 \pm 13.6$ & $165.1 \pm 4.4$  & $161.0 \pm 5.0$  \\
16 & 1  & $288.9 \pm 17.4$ & $219.4 \pm 8.9$  & $214.5 \pm 8.9$  \\
16 & 2  & $401.0 \pm 33.8$ & $279.9 \pm 13.1$ & $283.2 \pm 14.8$ \\
16 & 4  & $500.5 \pm 68.1$ & $325.6 \pm 16.7$ & $313.0 \pm 19.7$ \\
16 & 8  & $476.5 \pm 34.2$ & $334.9 \pm 16.7$ & $320.1 \pm 17.0$ \\
32 & 4  & $486.8 \pm 45.5$ & $463.1 \pm 42.0$ & $447.5 \pm 39.4$ \\
32 & 8  & $455.1 \pm 38.9$ & $483.5 \pm 44.0$ & $463.4 \pm 40.0$ \\
32 & 16 & $448.5 \pm 36.4$ & $308.2 \pm 18.4$ & $302.7 \pm 16.9$ \\
32 & 32 & $190.1 \pm 7.8$  & $140.5 \pm 3.9$  & $136.9 \pm 3.9$  \\
64 & 16 & $629.7 \pm 78.0$ & $495.8 \pm 46.3$ & $456.3 \pm 40.1$ \\
64 & 32 & $485.0 \pm 47.2$ & $201.8 \pm 8.8$  & $195.2 \pm 8.2$  \\
64 & 64 & $142.7 \pm 6.0$  & $92.3 \pm 2.0$   & $91.9 \pm 2.2$   \\
\bottomrule
\end{tabular}
\end{table}

Сравнение вариантов~3 и~10 на одинаковом размере задачи показывает, что
проверки границ различимо снижают производительность в 11 из 13
комбинаций параметров: вариант~10 медленнее варианта~3 на 21--58~\%
(наибольшее снижение~--- при $TS=64$, $WPT=32$). Различимого снижения
нет только при $TS=32$ и $WPT \in \{4, 8\}$. При лучших для каждого
варианта параметрах вариант~10 (495.8~GFLOPS при $TS=64$, $WPT=16$)
различимо медленнее варианта~3 (629.7~GFLOPS) примерно на 21~\%.
Сравнение двух размеров задачи для варианта~10 показывает, что при
всех 13 комбинациях результаты с неполными блоками и без них не
различимы. Следовательно, снижение производительности варианта~10
вызвано проверками границ, а не размерами матриц, и его величина
зависит от параметров. На задаче с неполными блоками лучшей для
варианта~10 оказалась комбинация $TS=32$, $WPT=8$~---
$463.4 \pm 40.0$~GFLOPS при среднем времени выполнения 0.290~с.
Причина того, что затраты на проверки границ зависят от параметров, в
работе не исследовалась.

\section{Ограничения исследования}

Реализованные варианты являются упрощёнными версиями шагов учебного
материала, а устройство вариантов~5, 7 и~10 отличается от кода
учебного материала, поэтому полученные результаты характеризуют именно
эти реализации. Шаги~8 и~9 не исследовались.

Перебор охватывает все значения параметров, допустимые при принятых
условиях (подраздел~\ref{sec:params}). Рабочие группы меньше
32 потоков, параметр \texttt{WIDTH} варианта~4 и параметр $WPT$
варианта~7 не проверялись, поэтому найденные параметры являются
лучшими только в пределах этих условий.

Каждая комбинация запускалась 10 раз без прогревочного запуска.
Разброс наиболее производительных комбинаций достиг
$\pm 123.3$~GFLOPS: например, у варианта~7 с $TS=64$, $WPT=4$ самый
медленный запуск показал 609.0~GFLOPS, а самый быстрый~--- 855.6~GFLOPS.
Из-за большого разброса часть различий между вариантами и между
комбинациями параметров оказалась неразличимой.

Результаты разных запусков скрипта могут различаться сильнее, чем
разброс внутри одного запуска: в предыдущем запуске перебора
(19.09.2026, результаты сохранены в истории репозитория) вариант~1 с
$TS=16$ показал $40.8 \pm 0.4$~GFLOPS, а в приведённых результатах~---
$52.3 \pm 0.5$~GFLOPS. Поэтому все сравнения в работе выполнены только
по результатам одного запуска скрипта.

Корректность проверялась выборочно, на 100 элементах матрицы результата
для каждой комбинации. Основное сравнение выполнено для одного размера
задачи ($4096 \times 4096 \times 4096$), а при проверке варианта~10
неполные блоки по измерению $M$ содержали одну строку.

Вклад каждого приёма оптимизации в производительность отдельно не
измерялся, а причины различий между вариантами и комбинациями
параметров не исследовались.

\section{Заключение}

В ходе работы получены следующие результаты:
\begin{itemize}
    \item на основе шагов учебного материала реализованы и запущены на
    графическом процессоре AMD Radeon 760M восемь вариантов SGEMM
    (варианты~1--7 и~10), корректность вычислений подтверждена
    выборочной проверкой для всех комбинаций параметров;
    \item при параметрах по умолчанию наибольшую среднюю
    производительность показал вариант~4~--- 496.5~GFLOPS, что
    примерно в 9.5 раза выше, чем у базового варианта, однако
    результаты вариантов~3, 4, 6 и~7 не различимы в пределах разброса;
    \item на основе ограничений устройства определены допустимые
    значения параметров, разработаны программа измерений и скрипт,
    перебирающий все 67 допустимых комбинаций;
    \item после подбора параметров наибольшую среднюю
    производительность показали варианты~6 и~7 с параметрами $TS=64$,
    $WPT=4$ (819.6 и 820.8~GFLOPS), что примерно в 15.7 раза выше,
    чем у базового варианта; их превосходство различимо над всеми
    вариантами, кроме варианта~3 с параметрами $TS=64$, $WPT=16$
    (629.7~GFLOPS);
    \item установлено, что поддержка матриц произвольного размера
    (вариант~10) при тех же параметрах различимо снижает
    производительность относительно варианта~3 в 11 из 13 комбинаций
    (на 21--58~\%), а наличие неполных блоков не влияет на
    производительность различимым образом;
    \item в репозитории настроена автоматическая проверка (CI), которая
    при каждом изменении кода собирает все варианты и проверяет
    корректность всех 67 комбинаций параметров.
\end{itemize}

\end{document}
